\begin{notechapter}
\chapter{Hierarchical Softmax and Huffman Trees}\label{ch:hierarchical-softmax}
\chaptersubtitle{A normalized distribution over $V$ words from $\log V$ binary decisions}

\begin{chapterintro}
The softmax over a vocabulary of $V$ words costs $O(dV)$ per prediction, because its
normalization sums over every word. Hierarchical softmax replaces the single choice among
$V$ words by a sequence of binary choices along a path in a binary tree whose leaves are
the words. Each binary choice is a logistic regression, the product of their probabilities
is automatically normalized, and only the nodes on one path are evaluated. With a Huffman
tree, frequent words get short paths and the average cost is about the entropy of the
word distribution, in bits. This chapter proves the normalization, derives the gradient,
builds a Huffman tree, and corrects several claims about how word2vec uses it. You will
learn
\begin{itemize}[nosep]
  \item the hierarchical softmax model and why it sums to $1$;
  \item its gradient: one logistic-regression gradient per node of the path;
  \item Huffman's algorithm, its optimality and the entropy bounds on its average path length;
  \item what word2vec actually does: a fixed tree built from counts, node vectors in place of
        output vectors.
\end{itemize}
\prerequisites{\cref{ch:word-embeddings}; the logistic sigmoid and binary cross-entropy
(\cref{ch:sigmoid-softmax}); entropy (\cref{ch:cross-entropy}).}
\end{chapterintro}

\section{The model}\label{hierarchical-softmax:sec:model}

Arrange the $V$ words as the leaves of a binary tree with $V - 1$ internal nodes; each internal node $n$
carries a vector $u_n\in\R^d$. For a word $w$, let $n_1(w),\dots,n_{L(w)}(w)$ be the internal nodes on the path from the root
to the leaf $w$, and let $d_k(w) = +1$ if the path goes to the left child at $n_k(w)$, $-1$ if it goes to the right. Given
a context vector $h$ (the centre-word vector $v_w$ in skip-gram, the averaged context in CBOW),
\[
  p(w\mid h) = \prod_{k=1}^{L(w)}\sigma\bigl(d_k(w)\,u_{n_k(w)}^\top h\bigr),
\]
where $\sigma$ is the logistic sigmoid. At each node, $\sigma(u_n^\top h)$ is the probability of going left and
$\sigma(-u_n^\top h) = 1 - \sigma(u_n^\top h)$ that of going right (Morin and Bengio, \emph{Hierarchical Probabilistic Neural
Network Language Model}, 2005).

\begin{proposition}[Normalization]\label{hierarchical-softmax:prop:normalized}
For every $h$ and every choice of node vectors, $\sum_{w}p(w\mid h) = 1$.
\end{proposition}

\begin{proof}
For a node $n$, let $m(n)$ be the sum of $p(w\mid h)$ over the leaves $w$ of the subtree rooted at $n$, each leaf's
product being restricted to the nodes from $n$ down. We show $m(n) = 1$ by induction from the
leaves up. For a leaf the product is empty, so $m = 1$. For an internal node with children $\ell$ and $r$,
every leaf below $\ell$ starts its product with $\sigma(u_n^\top h)$ and every leaf below $r$ with $\sigma(-u_n^\top h)$, so
$m(n) = \sigma(u_n^\top h)\,m(\ell) + \sigma(-u_n^\top h)\,m(r) = \sigma(u_n^\top h) + \sigma(-u_n^\top h) = 1$. At the root, $m$ is the required sum.
\end{proof}

No normalization over the vocabulary is ever computed: the normalization is built into the
tree. Computing $p(w\mid h)$ costs $L(w)$ dot products instead of $V$.

\section{The gradient}\label{hierarchical-softmax:sec:gradient}

The loss of one prediction is $\ell = -\log p(w\mid h) = -\sum_k\log\sigma(d_ks_k)$, with $s_k = u_{n_k}^\top h$. Write $t_k = (1 + d_k)/2$, which
is $1$ if the path goes left at $n_k$ and $0$ if it goes right. Each term is the binary cross-entropy of a logistic
regression with label $t_k$, so (\cref{backpropagation:sec:output})
\[
  \frac{\partial\ell}{\partial s_k} = \sigma(s_k) - t_k,\qquad
  \nabla_{u_{n_k}}\ell = \bigl(\sigma(s_k) - t_k\bigr)h,\qquad
  \nabla_h\ell = \sum_{k=1}^{L(w)}\bigl(\sigma(s_k) - t_k\bigr)u_{n_k} .
\]
Only the $L(w)$ node vectors of the path receive a gradient, and the context vector receives the
sum of their contributions.

\begin{example}[Four words]\label{hierarchical-softmax:ex:four}
Take the tree with root $R$, whose left child $A$ has the leaves cat (left) and dog (right), and whose right
child $B$ has the leaves apple (left) and banana (right). With $h = (1, -0.5)$, $u_R = (0.4, 0.2)$, $u_A = (1, 1)$, $u_B = (-0.5, 0.8)$,
the scores are $s_R = 0.3$, $s_A = 0.5$, $s_B = -0.9$, and $\sigma(s_R) = 0.5744$, $\sigma(s_A) = 0.6225$, $\sigma(s_B) = 0.2891$. Then
\begin{align*}
  p(\text{cat}) &= 0.5744\times0.6225 = 0.3576, & p(\text{dog}) &= 0.5744\times0.3775 = 0.2169,\\
  p(\text{apple}) &= 0.4256\times0.2891 = 0.1230, & p(\text{banana}) &= 0.4256\times0.7109 = 0.3025,
\end{align*}
which sum to $1$. For the target ``dog'' (left at $R$, right at $A$), $\ell = -\log0.2169 = 1.5284$, and
$\partial\ell/\partial s_R = 0.5744 - 1 = -0.4256$, $\partial\ell/\partial s_A = 0.6225 - 0 = 0.6225$, so
\begin{gather*}
  \nabla_{u_R}\ell = (-0.4256,\ 0.2128),\qquad\nabla_{u_A}\ell = (0.6225,\ -0.3112),\qquad\nabla_{u_B}\ell = 0,\\
  \nabla_h\ell = -0.4256\,u_R + 0.6225\,u_A = (0.4522,\ 0.5373) .
\end{gather*}
A gradient step increases $s_R$ (more mass to the left subtree) and decreases $s_A$ (more mass to dog
within it), and leaves $u_B$ unchanged.
\end{example}

\section{Huffman trees}\label{hierarchical-softmax:sec:huffman}

Any binary tree gives a valid model, and the cost of a prediction is the depth of the target
word. Training visits each word in proportion to its frequency $f_w$, so the expected cost per
example is the average depth $\bar L = \sum_wf_wL(w)$. The tree that minimizes it is the Huffman tree, built
for optimal prefix codes (Huffman, \emph{A Method for the Construction of Minimum-Redundancy
Codes}, 1952):
\begin{enumerate}[nosep]
  \item create one leaf per word, with weight its count;
  \item repeatedly remove the two trees of smallest weights and join them under a new
        internal node whose weight is the sum of theirs;
  \item stop when one tree is left.
\end{enumerate}
Ties can be broken arbitrarily: different choices give different trees with the same $\bar L$.

\begin{theorem}[Huffman]\label{hierarchical-softmax:thm:huffman}
Among all binary trees with the words as leaves, the Huffman tree minimizes $\bar L = \sum_wf_wL(w)$, and
\[
  H(f)\leq\bar L < H(f) + 1,\qquad H(f) = -\sum_wf_w\log_2f_w .
\]
\end{theorem}

Optimality is proved by an exchange argument, and the bounds by Kraft's inequality; see Cover
and Thomas, \emph{Elements of Information Theory}, 2006, Chapter~5. A balanced tree has depth
$\lceil\log_2V\rceil$ for every word, and $H(f)\leq\log_2V$: for the skewed frequencies of natural language, the
Huffman tree is shorter on average. The worst case, however, can be much longer than $\log_2V$
(up to $V - 1$ for a word of very low frequency), so the claim that hierarchical softmax costs
$O(\log V)$ is a statement about the average cost.

\begin{figure}[ht]
\centering
\begin{forest}
for tree={circle, draw=bookaccent, thick, fill=bookaccent!8, font=\scriptsize, minimum size=3.2em, inner sep=1pt,
  s sep=12pt, l sep=6pt, align=center},
where n children=0{fill=bookexample!20, draw=bookexample!80!black}{}
[39
  [17
    [dog\\8]
    [9
      [banana\\5]
      [elephant\\4]
    ]
  ]
  [22
    [cat\\10]
    [12
      [apple\\6]
      [6
        [zebra\\3]
        [3
          [giraffe\\2]
          [lion\\1]
        ]
      ]
    ]
  ]
]
\end{forest}
\caption{A Huffman tree for the counts cat $10$, dog $8$, apple $6$, banana $5$, elephant $4$, zebra $3$, giraffe $2$, lion $1$. The
merges are $1 + 2$, $3 + 3$, $4 + 5$, $6 + 6$, $8 + 9$, $10 + 12$, $17 + 22$; internal nodes show their weights.}
\label{hierarchical-softmax:fig:huffman}
\end{figure}

\begin{example}[Average path length]\label{hierarchical-softmax:ex:huffman}
In \cref{hierarchical-softmax:fig:huffman}, the depths are $2$ for cat and dog, $3$ for apple, banana and elephant, $4$ for zebra and
$5$ for giraffe and lion. With $f_w = \text{count}/39$, $\bar L = (2\cdot18 + 3\cdot15 + 4\cdot3 + 5\cdot3)/39 = 108/39 = 2.769$, against $3$ for the
balanced tree with $8$ leaves, and the entropy is $H = 2.745$ bits: $H\leq\bar L < H + 1$, and here $\bar L$ is within $0.03$ bit
of the bound. The rarest words pay with a depth of $5$.
\end{example}

\section{Hierarchical softmax in word2vec}\label{hierarchical-softmax:sec:word2vec}

In word2vec with hierarchical softmax (Mikolov et al., 2013), the probability of a context word
$c$ given the centre word $w$ is $p(c\mid v_w)$ computed by the tree, with $h = v_w$. Three points are often stated
incorrectly.
\begin{itemize}
  \item \textbf{The tree is fixed.} It is built once, by Huffman's algorithm on the word counts of the
        training corpus, before training starts, and it never changes. Only the vectors are
        learned. Methods that learn the tree exist, but as a separate step: Mnih and Hinton
        (\emph{A Scalable Hierarchical Distributed Language Model}, 2009) build it by clustering
        word vectors learned beforehand.
  \item \textbf{The node vectors replace the output vectors.} There are $V - 1$ node vectors $u_n$ and no output
        vector per word; the logistic classifier at a node \emph{is} its vector (word2vec uses no
        bias), so ``updating the node vectors and the binary classifiers'' describes one
        update, not two. The input vectors $v_w$ are the embeddings.
  \item \textbf{GloVe does not use it.} GloVe has no softmax to approximate
        (\cref{word-embeddings:sec:glove}). FastText does offer hierarchical softmax, mainly for
        its supervised text classifier with many labels.
\end{itemize}

\begin{pitfall}[What the Huffman tree optimizes]
The Huffman tree minimizes the \emph{expected cost per training example}, because training
visits frequent words more often. It does not make the model more accurate, and it does not
group similar words: the leaves are placed by frequency alone, so ``cat'' and ``dog'' may be
far apart while ``cat'' and ``of'' share a long prefix. Each internal node then has to
separate two arbitrary groups of words with one linear classifier, which is a real
constraint on the model; trees built from word similarity (as in Mnih and Hinton) make these
binary decisions easier.
\end{pitfall}

\begin{chaptersummary}
\begin{itemize}
  \item Hierarchical softmax writes $p(w\mid h) = \prod_k\sigma(d_ku_{n_k}^\top h)$ along the path from the root to the leaf $w$; it is
        normalized for any parameters because $\sigma(s) + \sigma(-s) = 1$ at every node.
  \item Each node is a logistic regression: $\nabla_{u_{n_k}}\ell = (\sigma(s_k) - t_k)h$, and only the path is updated.
  \item The Huffman tree minimizes the average depth $\bar L$, with $H\leq\bar L < H + 1$; rare words can be deep.
  \item In word2vec the Huffman tree is built once from counts and fixed; its $V - 1$ node vectors replace
        the output vectors; GloVe does not use hierarchical softmax.
\end{itemize}
\end{chaptersummary}

\section*{Exercises}
\addcontentsline{toc}{section}{Exercises}

\begin{exercise}\label{hierarchical-softmax:exr:1}
In \cref{hierarchical-softmax:ex:four}, compute the gradients for the target ``apple'', and check that
$\nabla_{u_A}\ell = 0$.
\solutionlink{hierarchical-softmax:exr:1}
\end{exercise}

\begin{exercise}\label{hierarchical-softmax:exr:2}
Show that a binary tree with $V$ leaves in which every internal node has two children has exactly $V - 1$
internal nodes.
\solutionlink{hierarchical-softmax:exr:2}
\end{exercise}

\begin{exercise}\label{hierarchical-softmax:exr:3}
Build the Huffman tree for the counts $a: 45$, $b: 13$, $c: 12$, $d: 16$, $e: 9$, $f: 5$, give the depth of each word and $\bar L$, and
compare with $H$ and with a balanced tree.
\solutionlink{hierarchical-softmax:exr:3}
\end{exercise}

\begin{exercise}\label{hierarchical-softmax:exr:4}
For $V$ words with counts $2^{V-2}, 2^{V-3},\dots,2, 1, 1$, show that the Huffman tree is a ``caterpillar'' in which the two
rarest words have depth $V - 1$, and compute $\bar L$ and $H$.
\solutionlink{hierarchical-softmax:exr:4}
\end{exercise}
\end{notechapter}
